% English Abstract

The continuous development of wireless communication and Internet of Things (IoT) technologies has driven the deployment of emerging applications such as virtual reality, augmented reality, and autonomous driving, which in turn impose stringent requirements on the computational capabilities at the network edge. Edge computing effectively reduces response latency by offloading computational tasks to edge nodes. However, with the surge of application data and the expansion of computational task scales, a single node can hardly support large-scale computational tasks independently, motivating the development of distributed edge computing. Distributed edge computing aggregates computational resources from multiple edge nodes, decomposing large-scale tasks into multiple subtasks for parallel processing, thereby improving computational efficiency. Nevertheless, due to the heterogeneity in hardware configurations, load states, and network environments of edge nodes, some nodes may become stragglers due to resource competition or sudden failures, with computational speeds far behind other nodes, causing the overall computational latency to be limited by the slowest node. To mitigate the straggler effect, coded edge computing has been proposed as an emerging paradigm integrating erasure coding theory with distributed computing. By introducing redundancy in the task distribution phase, the central node only needs to collect computation results from a subset of faster nodes to reconstruct the original task results. However, most existing research on coded edge computing focuses on computational latency optimization, overlooking the communication latency that may arise from channel fading during the downlink transmission phase. In fact, the redundancy introduced by coded computing during task distribution creates strong correlation among computation results from different edge nodes. If this correlation can be exploited to compress the computation results, the transmission rate required by each node can be reduced, thereby alleviating communication latency. Distributed Source Coding (DSC), as an efficient compression technique designed for correlated sources, provides a theoretical foundation for optimizing downlink transmission in coded edge computing. Addressing these issues, this thesis investigates the joint optimization of coded computation result transmission and DSC, targeting two scenarios: complete recovery and linear recovery of computation results. The main research contents and contributions are as follows:

(1) For coded computing scenarios requiring complete recovery of subtask computation results, a DSC-assisted coded computation result transmission scheme for complete recovery is proposed. This scheme leverages the linear correlation among coded computation results and combines the Slepian-Wolf (SW) theorem to characterize the theoretical compression region for achieving complete recovery of computation results without inter-node communication. Since the compression limit given by the SW theorem lacks directly implementable coding schemes, a coding scheme based on Low-Density Parity-Check (LDPC) code with syndrome puncturing is further designed, enabling edge nodes to adaptively adjust compression rates according to channel states, achieving flexible reallocation of transmission loads among nodes, and extending the scheme to high-dimensional data scenarios. Based on the proposed scheme, a downlink transmission latency minimization model is established. For this non-convex optimization problem, a joint optimization algorithm based on Convex-Concave Procedure (CCCP) is designed to jointly optimize source compression rates and transmit power of each edge node. Simulation results demonstrate that the proposed scheme significantly reduces downlink transmission latency, especially in scenarios with poor channel conditions or limited transmit power, and exhibits good robustness to variations in the number of straggler nodes.

(2) For coded computing scenarios requiring only the linear function value of computation results (such as gradient aggregation in distributed machine learning), since complete recovery of individual subtask results is unnecessary, there exists room for further compression. Accordingly, a DSC-assisted coded computation result transmission scheme for linear recovery is proposed. This scheme combines nested linear codes to characterize the achievable rate region for scenarios requiring only the linear function value of computation results. It is proved that the achievable rate region for linear recovery is strictly larger than the classical SW region, thereby achieving higher compression efficiency than complete recovery. Based on the proposed scheme, an inner bound of the achievable rate region is provided, namely the union of the SW boundary and the linear recovery boundary, which reduces the complexity of solving the optimization problem. Furthermore, a corresponding latency minimization model is established, and an optimization algorithm based on CCCP is designed. Simulation results demonstrate that the downlink transmission performance of the proposed scheme outperforms both independent coding schemes and SW coding schemes, with more pronounced advantages in scenarios with many straggler nodes.

In summary, by exploiting the strong correlation among tasks in coded edge computing, this thesis combines DSC with coded computing, providing a new perspective and effective solutions for addressing communication bottlenecks in edge computing systems.